% ============================================================ real scans
\section{Reconstruction from real scans}\label{sec:scans}
Every comparison so far has used clouds sampled from clean synthetic meshes with \emph{mesh-derived} normals.
Real $3$D scans are harder in three ways at once, and all three are exactly what a reconstruction method must
survive in the wild: they carry \emph{sensor noise} (each point is jittered off the true surface), they are
\emph{incomplete} (self-occlusion and limited views leave holes and one-sided coverage), and --- crucially ---
they have \emph{no normals at all}.  We therefore feed every method the same honest input: the raw scan
vertices, centred and scaled into the $[-1,1]$ frame and farthest-point subsampled to $4096$ points, with
per-point normals \emph{estimated} from the points themselves (a local PCA fit, oriented for global
consistency by a minimum-spanning-tree propagation over the tangent planes, via Open3D~\cite{open3d}).  These
estimated normals are noisier and occasionally mis-oriented where the surface is thin or sparsely sampled ---
precisely the regime our analytic gate (\S\ref{sec:mixed}) is most sensitive to --- so this is a genuine
stress test of the whole pipeline, anchor (\S\ref{sec:splat}) and gate included, not just of the network.

\paragraph{No ground truth: a qualitative \emph{and} validity evaluation.}  A real scan has no reference
mesh, so the F-score, Chamfer and signed-distance metrics of \S\ref{sec:baselines} are simply
\emph{undefined} here --- there is nothing to compare against.  We do not invent a pseudo-ground-truth (e.g.\
a Poisson mesh of the scan), which would only measure agreement with one particular baseline.  Instead the
evaluation is twofold and both halves are GT-free: \emph{(i)} \textbf{qualitative} --- does the reconstructed
surface look like the object, and does it degrade gracefully under real noise and occlusion rather than
collapsing? --- and \emph{(ii)} \textbf{mesh validity}, the two defect counts that need no reference: the
number of connected \textbf{components} (a clean object should be few; a fragment soup is many) and the number
of \textbf{self-intersecting faces} (\textsc{self-X}), both read off the output mesh by the same
\texttt{mesh\_defects} routine used in Table~\ref{tab:quant}.  These are reported under each panel of
Figure~\ref{fig:cmp_scans}.

\begin{figure*}[t]
\centering\includegraphics[width=\textwidth]{figs/cmp_scans.png}
\caption{\textbf{Reconstruction from real scans} (no ground truth; $4096$ points, \emph{estimated} normals).
Rows are scanned objects; columns are the input point cloud, \textsc{spsr}~\cite{kazhdan2013} and
\textsc{bpa}~\cite{bernardini1999} (Open3D~\cite{open3d}), \textsc{apss}~\cite{apss} and
\textsc{rimls}~\cite{rimls} (PyMeshLab~\cite{pymeshlab}), and \emph{ours} (the unified mixed-base model,
\S\ref{sec:mixed}, queried at $R{=}128$).  Because real scans have \emph{no} reference surface there is no
F-score or Chamfer; instead each reconstructed panel prints two GT-free mesh-validity numbers --- the
connected-component count and the self-intersecting-face count (\textsc{self-X}).  Normals are \emph{estimated}
from the points (local PCA, globally oriented), not read from a mesh, so this also stresses the per-point gate,
which depends on neighbouring normal directions.}
\label{fig:cmp_scans}
\end{figure*}

\paragraph{What the scans show.}  The qualitative picture is consistent with the synthetic comparison of
\S\ref{sec:baselines}, and the failure modes are the expected ones.  \textsc{spsr}, a watertight \emph{signed}
Poisson solver, is the most robust to scan noise --- its global smoothing closes holes and yields a single
clean component --- but it pays for that on \emph{open} or partially-scanned geometry, where it hallucinates a
closed bubble over the missing back side (the over-fill of \S\ref{sec:limits}, now from real occlusion rather
than synthetic thinness).  \textsc{bpa} hugs the measured points faithfully and so reproduces fine detail, but
on noisy, unevenly-sampled scans it leaves boundary gaps and shatters into many components.  \textsc{apss} and
\textsc{rimls}, the moving-least-squares surfaces, smooth the noise into a single connected blob yet carry
heavy self-intersection (the \textsc{self-X} counts mirror Table~\ref{tab:quant}).  Ours reconstructs the
overall body and survives the estimated, occasionally-flipped normals, routing closed regions through the
signed base and thin/occluded regions through the unsigned band so it does not over-fill the way \textsc{spsr}
does; its weaknesses are equally those of \S\ref{sec:baselines} --- residual floating fragments raise its
component count, and where the estimated normals are inconsistent on a thin or sparsely-covered patch the gate
can misfire, the failure case anticipated in \S\ref{sec:limits-gate}.  The honest reading is the same as on
synthetic data: no method is clean on real scans, and ours is competitive --- avoiding the worst over-fill and
self-intersection pathologies --- rather than dominant.
